\unnumchapter{Введение}

Современные сервисы всё чаще работают не в одном, а сразу в нескольких
Kubernetes-кластерах~\cite{k8s}: в разных регионах, у разных провайдеров, с
разными требованиями к хранению и обработке данных. Такой набор кластеров под
общим управлением называют федерацией. Два свойства федерации определяют
круг задач, которые приходится решать при её эксплуатации. Во-первых,
кластеров несколько, поэтому отказ одного из них является штатным событием, и
вопрос состоит не в том, откажет ли площадка, а в том, что в этот момент
произойдёт с сервисом. Во-вторых, размещать нагрузку на произвольной площадке
нельзя: регион, сертификации, требования к безопасности и доступные ресурсы
сужают множество допустимых кластеров для каждой конкретной нагрузки.

Для каждой из подзадач эксплуатации федерации существуют зрелые инструменты:
системы мониторинга, планировщики, движки политик, средства резервного
копирования. Однако эти инструменты не связаны между собой. Система
мониторинга обнаруживает, что кластер деградирует, но планировщик об этом не
знает и продолжает направлять туда новую нагрузку. Политики размещения
проверяются в границах одного кластера и не учитывают свойства соседних.
Процедура восстановления после аварии описана текстом в регламенте и
выполняется человеком. В результате время реакции на инцидент определяется
квалификацией и доступностью дежурного инженера, а не свойствами самой
системы.

В настоящей работе рассматривается подход, при котором перечисленные
звенья объединяются в сквозной контур управления: наблюдаемость показывает
фактическое качество работы каждого кластера, целевые показатели уровня
обслуживания (SLO) переводят это качество в измеримый критерий, политики
отсекают недопустимые площадки, планировщик уровня федерации выбирает
лучшую из оставшихся с учётом их текущего состояния, а контроллер
аварийного восстановления (DR) воспроизводимо возвращает систему в рабочее
состояние. Главное отличие такого контура от набора отдельных инструментов
состоит в том, что решение о размещении принимается с учётом текущего
состояния кластеров, а не только их статических атрибутов.

Принципиальным для работы является требование измеримости: каждый механизм
контура оценивается экспериментально, путём сравнения с вариантом без него на
одном и том же воспроизводимо внесённом отказе. Механизм, эффект которого не
измерен, ценности для исследования не имеет.

В работе исследуются следующие вопросы. Как определить и измерить показатели
надёжности мультикластерного сервиса, агрегируя метрики нескольких
независимых кластеров? Какие механизмы федеративного контура~--- переключение
трафика при отказе и размещение нагрузки с учётом метрик~--- и в какой мере
улучшают показатели надёжности при типовых сбоях: отказе кластера и
деградации сети? Как обеспечить воспроизводимый цикл аварийного
восстановления, от обнаружения отказа до управляемого возврата трафика, с
фиксируемым временем восстановления и числом потерянных запросов? Как
формализовать политики размещения по регионам, сертификациям и ресурсам и
обеспечить их исполнение без ручного вмешательства? Наконец, центральный
вопрос работы: можно ли построить планировщик уровня федерации, учитывающий
текущие метрики качества кластеров и политики размещения, и насколько он
превосходит наивный выбор кластера? Остальные вопросы создают данные и
условия, при которых на центральный вопрос можно ответить количественно.

Существующие федеративные планировщики выбирают кластер по статическим
атрибутам: меткам и ресурсной ёмкости. В предлагаемом подходе в выборе
участвуют текущие доля ошибок и задержка, а политики, наблюдаемость и
восстановление связаны через общий каталог кластеров. Новизна работы
заключается не в самом алгоритме выбора, а в замкнутом контуре и в том, что
эффект каждого его звена измерен относительно базового варианта на
одинаковых сценариях отказов.

Практическая значимость работы определяется тем, что все компоненты
контура, сценарии экспериментов и сырые данные прогонов воспроизводимы:
стенд и каждый эксперимент разворачиваются из репозитория одной командой,
план аварийного восстановления является исполняемым, а не текстовым, а
методика измерения переносима на реальные мультикластерные системы.
